% Work-MAT | 2. rocnik | pripravovana hodina 40
% Návaznost: uživatelské připravy.zip (hodiny 01–20), autorská učebnice na GitHubu.
% Realisticky.cz: didaktická inspirace v posloupnosti a typech úloh; zadání jsou vlastní.
\documentclass[11pt,a4paper]{article}
\input{../../sablony/pripravy/priprava.sty}
\nastavPripravu{2. ročník}{Logaritmické rovnice}{40}{Rovnice s kvadratickým argumentem}
\begin{document}
\blok{Motivace a návaznost}
V logaritmu se může objevit kvadratický výraz. Na začátku připomeň, jak řešit nerovnici $x^2-1>0$ a jak rozhodovat, zda kořen patří do definičního oboru.
\poznamka{Orientační nabídka pro 45--55 minut: krátké zopakování a motivace, nový výklad, tři společné vzorové úlohy, samostatná práce. Vyber podle tempa třídy; zbývající zadání slouží jako rezerva.}
\blok{Výklad}
Podmínka $F(x)>0$ pro logaritmus se zjišťuje dříve než úprava rovnice; případné další algebraické kořeny se nepřebírají automaticky.
\blok{Společné příklady}
\begin{prikladReseniBox}[title={Příklad 1}]
Vyřeš, resp. urči: $\log_3(x^2-1)=2$.
\tcblower
\[\begin{aligned}x^2-1>0 \\ x^2-1=9 \\ x=\pm\sqrt{10}\end{aligned}\]
\end{prikladReseniBox}
\begin{prikladReseniBox}[title={Příklad 2}]
Vyřeš, resp. urči: $\log_2(x^2-3)=2$.
\tcblower
\[\begin{aligned}x^2>3 \\ x^2=7 \\ x=\pm\sqrt7\end{aligned}\]
\end{prikladReseniBox}
\poznamka{Zkoušku nevynechávej: součástí dnešní hodiny je i řešení nerovnice určující definiční obor.}
\blok{Samostatné procvičování}
\begin{zadaniReseniBox}[title={Samostatná práce 1}]
\begin{enumerate}[label=\alph*)]
\item $\log_{2}\left(x^2-1\right)=1$.
\item $\log_{3}\left(x^2-2\right)=2$.
\item $\log_{5}\left(x^2-3\right)=3$.
\item $\log_{2}\left(x^2-1\right)=1$.
\item $\log_{3}\left(x^2-2\right)=2$.
\item $\log_{5}\left(x^2-3\right)=3$.
\end{enumerate}
\tcblower
\begin{enumerate}[label=\alph*)]
\item $x=\pm\sqrt{3}$.
\item $x=\pm\sqrt{11}$.
\item $x=\pm\sqrt{128}$.
\item $x=\pm\sqrt{3}$.
\item $x=\pm\sqrt{11}$.
\item $x=\pm\sqrt{128}$.
\end{enumerate}
\end{zadaniReseniBox}
\blok{Dobrovolné rozšíření}
\begin{zadaniReseniBox}[title={Rozšíření 1}]
\begin{enumerate}[label=\alph*)]
\item $Vytvoř vlastní úlohu tohoto typu a se spolužákem porovnej podmínky a výsledek.$.
\end{enumerate}
\tcblower
\begin{enumerate}[label=\alph*)]
\item Zkontroluj definice a argumenty..
\end{enumerate}
\end{zadaniReseniBox}
\blok{Úkol na doma}
U každého logaritmu vypiš podmínky argumentu a proveď zkoušku. U rovnic vyznač případně nepřípustné kořeny, u nerovnic odpověz intervalem.
\begin{zadaniReseniBox}[title={Domácí úkol}]
\begin{enumerate}[label=\alph*)]
\item $\log_{2}\left(x^2-1\right)=1$.
\item $\log_{3}\left(x^2-2\right)=2$.
\item $\log_{5}\left(x^2-3\right)=3$.
\item $\log_{2}\left(x^2-1\right)=1$.
\item $\log_{3}\left(x^2-2\right)=2$.
\item $\log_{5}\left(x^2-3\right)=3$.
\item $\log_{2}\left(x^2-1\right)=1$.
\item $\log_{3}\left(x^2-2\right)=2$.
\end{enumerate}
\tcblower
\begin{enumerate}[label=\alph*)]
\item $x=\pm\sqrt{3}$.
\item $x=\pm\sqrt{11}$.
\item $x=\pm\sqrt{128}$.
\item $x=\pm\sqrt{3}$.
\item $x=\pm\sqrt{11}$.
\item $x=\pm\sqrt{128}$.
\item $x=\pm\sqrt{3}$.
\item $x=\pm\sqrt{11}$.
\end{enumerate}
\end{zadaniReseniBox}
\poznamka{Domácí úkol obsahuje 8 položek; předpokládaná domácí práce přibližně 15--30 minut. Po větších blocích jsou vědomě promíchané dříve probrané typy.}
\end{document}
